% Scope: Develop only the mathematics and electromagnetism needed for spin dynamics, encoding, noise and estimation; derive prerequisites at their first useful level.
% Status: architecture only; no chapter prose or completed problems yet.
% Prerequisites: None; conventions in the front matter.
% Planned worked examples: Phasor signs, sampling limits and covariance propagation.
% Planned figures: Complex plane, convolution, sampling replicas.
% Planned challenge problems: Prove an aliasing ambiguity; optimize a noisy linear estimator.
\chapter{Mathematical and Physical Foundations}
\label{ch:01}

\section{Complex representation and dimensional reasoning}
\label{sec:01:complex-representation-and-dimensional-reasoning}
\subsection{Complex numbers, Euler's formula and phasors}
\subsection{Frequency, angular frequency and phase}
\subsection{SI units, scaling and dimensional checks}

\section{Linear algebra for spin and encoding}
\label{sec:01:linear-algebra-for-spin-and-encoding}
\subsection{Vector spaces, inner products and adjoints}
\subsection{Eigenvalues, eigenvectors and matrix exponentials}
\subsection{Singular values, rank and conditioning}

\section{Fourier analysis and linear systems}
\label{sec:01:fourier-analysis-and-linear-systems}
\subsection{Transform pairs and distributions}
\subsection{Convolution, filtering and impulse responses}
\subsection{Sampling, Nyquist limits and discrete transforms}

\section{Probability and measurement noise}
\label{sec:01:probability-and-measurement-noise}
\subsection{Random variables, expectations and covariance}
\subsection{Gaussian complex noise and magnitude statistics}
\subsection{Likelihood, bias and propagation of uncertainty}

\section{Estimation and experimental information}
\label{sec:01:estimation-and-experimental-information}
\subsection{Least squares and maximum likelihood}
\subsection{Fisher information and the Cramer--Rao bound}
\subsection{Identifiability and nuisance parameters}

\section{Fields, dipoles and angular momentum}
\label{sec:01:fields-dipoles-and-angular-momentum}
\subsection{Maxwell equations and quasistatic limits}
\subsection{Magnetic dipole energy and torque}
\subsection{Angular momentum and rotating coordinates}

\section{Worked examples}
\label{sec:01:worked-examples}
% Planned calculations: Phasor signs, sampling limits and covariance propagation.
\subsection{Derivation and quantitative calculation}
\subsection{Assumptions, limiting cases and scanner interpretation}

\section{Challenge problems}
\label{sec:01:challenge-problems}
% Problem design: Prove an aliasing ambiguity; optimize a noisy linear estimator.
% Use challengeproblem and stable labels prob:01:<topic>.
% Complete solutions belong in Appendix E, section for this chapter.
% Use \begin{solution}{prob:01:<topic>} ... \end{solution} there.
\subsection{Derivation and quantitative design}
\subsection{Model criticism and integrated reasoning}
